% Interfacial equilibrium, drawn generic: no solubility value and no numbered axes, since
% neither has been introduced at this point of the document. The slope of the Henry line is
% arbitrary and must stay so.
% Requires in the preamble: \usepackage{tikz}
%   \usetikzlibrary{positioning, fit, backgrounds, arrows.meta}
\begin{tikzpicture}[
    >=Stealth,
    axis/.style={->, thick, black!55},
    henry/.style={very thick, black!80},
    guide/.style={dashed, black!45, line width=0.5pt},
    flux/.style={->, thick, violet!70!black},
    tend/.style={->, line width=1.1pt, orange!70!black},
    font=\small,
]
% the Henry line separates the plane in two: above it the liquid holds more tritium than the
% bubble can balance, below it the bubble does
\fill[green!8] (0,0) -- (6.95,3.128) -- (6.95,4.5) -- (0,4.5) -- cycle;

\draw[axis] (0,0) -- (6.95,0) node[below left=1pt and -6pt] {$\PT$ in the bubble};
\draw[axis] (0,0) -- (0,4.5) node[above right=-3pt and -18pt] {$\CT$ in the liquid};
\draw[henry] (0,0) -- (6.5,2.925)
    node[sloped, above, pos=0.74, inner sep=3pt, black] {$\CT=K_s\,\PT$};

% state of the liquid and of the bubble it surrounds, taken above the line so that tritium
% crosses the interface outwards, which is the extraction direction
\draw[guide] (0,3.3) -- (1.6,3.3);
\draw[guide] (1.6,3.3) -- (1.6,0);
\filldraw[black!80] (1.6,3.3) circle (1.5pt);
\filldraw[draw=black!80, fill=white, line width=0.7pt] (1.6,0.72) circle (1.5pt);
\node[anchor=south east, inner sep=3pt] at (1.6,3.3) {bulk state};
\node[anchor=north west, inner sep=3pt] at (1.6,0.72) {interface, at equilibrium};

% the interfacial flux, proportional to the vertical gap to the Henry line
\draw[flux] (1.6,3.3) -- (1.6,0.80);
\node[rotate=90, anchor=south, violet!70!black, inner sep=4pt] at (1.6,2.05) {\JT};

% where the state travels: the liquid gives up tritium and the bubble takes it, so the point
% moves down and to the right until it reaches the line
\draw[tend] (2.05,3.16) .. controls (3.0,2.95) and (3.3,2.70) .. (3.90,2.45);
\node[anchor=south west, orange!70!black, font=\small\itshape, inner sep=1pt]
    at (2.30,3.10) {tends to equilibrium};

\node[anchor=north east, black!55, font=\footnotesize\itshape] at (6.95,4.5)
    {liquid supersaturated};
\node[anchor=east, black!55, font=\footnotesize\itshape] at (6.7,1.15)
    {bubble supersaturated};
\end{tikzpicture}
